\chapter{Un socle pour l'IA en hydrologie}
\label{ch:socle}

\questionchapitre{Ces développements survivront-ils au projet qui les a financés, et à quelles
conditions peuvent-ils fonder un environnement d'outils plus vaste ?}

\section{Pourquoi la plupart des développements de recherche ne se réutilisent pas}

La réponse par défaut est négative. Un logiciel de recherche est ordinairement écrit pour produire un résultat, puis abandonné : il
suppose un environnement particulier, mêle le traitement et son affichage, dépend de fichiers
qui ne sont plus là, et sa reprise coûte plus cher que sa réécriture.

Affirmer qu'un développement constitue un « socle » n'a donc de valeur que si l'on peut désigner
les propriétés qui le rendent réutilisable et montrer qu'elles sont effectivement acquises.

\section{Les propriétés acquises}

La couche analytique de l'entrepôt se consomme en SQL, et rien d'autre. Un utilisateur n'a besoin
de connaître ni l'outil d'ingestion, ni la bibliothèque de transformation, ni l'orchestrateur :
un carnet de calcul, un tableur connecté ou une application tierce s'y branchent également. La
conséquence est structurante : \textbf{l'amont peut être remplacé sans que l'aval s'en
aperçoive}. Un outil d'ingestion devenu obsolète se remplace sans
toucher aux consommateurs. Peu de logiciels de recherche offrent cette propriété, et c'est
elle qui distingue une infrastructure d'un script.

Le cœur de calcul de la plateforme, lui, ne dépend d'aucune interface. La contrainte a été posée
dès l'origine, au chapitre~\ref{ch:observatoire}, et son effet sur la valeur résiduelle des
travaux est direct : \textbf{même si l'application web était abandonnée, les méthodes
resteraient utilisables}. La contribution n'est pas prisonnière de son emballage.

Le reste a été acquis en chemin. Des indicateurs calculés une fois pour tous les consommateurs,
sans second calcul susceptible de diverger, et gardés par des tests qui échouent si les deux
dépôts cessent de s'accorder (chapitre~\ref{ch:indicateurs}). Une pile conteneurisée dont les
volumes de données survivent à l'arrêt des conteneurs, et dont la restauration a été vérifiée à
l'échelle réelle (chapitre~\ref{ch:entrepot}). Une documentation refondue en fin de période pour
la reprise, jusqu'aux pièges qui produisent des résultats faux sans émettre d'erreur
(chapitre~\ref{ch:bilan}).

\section{Le patron généralisable}

Au-delà des deux dépôts, ces travaux instancient un enchaînement (figure~\ref{fig:patron}) qui
n'a rien de spécifique aux eaux souterraines.

\begin{figure}[H]
\centering
\begin{tikzpicture}[
  font=\sffamily\small,
  et/.style={draw=black!70, fill=black!3, minimum width=2.6cm, minimum height=1.3cm,
             align=center, font=\sffamily\scriptsize},
  fl/.style={-{Stealth[length=2.4mm]}, draw=black!70, thick},
  note/.style={font=\sffamily\scriptsize\itshape, color=black!55, align=center,
               text width=2.55cm},
]
\node[et] (e1) {\textbf{Sources}\\ouvertes,\\hétérogènes};
\node[et, right=0.7cm of e1] (e2) {\textbf{Entrepôt}\\consolidé,\\interrogeable};
\node[et, right=0.7cm of e2] (e3) {\textbf{Indicateurs}\\normalisés\\et validés};
\node[et, right=0.7cm of e3] (e4) {\textbf{Modèles}\\métier et\\appris};
\node[et, right=0.7cm of e4, very thick, draw=black] (e5) {\textbf{Explication}\\et\\scénarisation};

\foreach \a/\b in {e1/e2, e2/e3, e3/e4, e4/e5} \draw[fl] (\a) -- (\b);

\node[note, below=0.32cm of e2] {coût payé une seule fois};
\node[note, below=0.32cm of e3] {comparabilité garantie ici};
\node[note, below=0.32cm of e5] {condition de l'usage en décision};
\end{tikzpicture}
\caption{L'enchaînement mis en \oe uvre. Chaque étape suppose la précédente ; escamoter l'une
d'elles produit des résultats incomparables ou inutilisables.}
\label{fig:patron}
\end{figure}

L'enseignement principal de l'année tient dans la lecture de gauche à droite de cette figure :
\textbf{on ne peut pas commencer par la droite}. C'est pourtant ce que fait le plus grand nombre
de projets d'intelligence artificielle appliquée à l'environnement, qui investissent dans le
modèle et découvrent ensuite que la donnée ne permet pas de conclure.

\section{Extensions, par coût croissant}

\begin{center}
\begin{tabularx}{\textwidth}{@{}lXl@{}}
\toprule
\textbf{Extension} & \textbf{Ce qu'elle suppose} & \textbf{Coût} \\
\midrule
Nouvelles stations, autre territoire & Rien : la couverture est déjà nationale et la chaîne est
paramétrée par configuration & nul \\
Nouveaux indicateurs standardisés & Un ajustement statistique et sa validation, selon le
protocole existant & faible \\
Nouvelles architectures de prévision & Une entrée au catalogue ; la fabrique de modèles et la
traçabilité sont génériques & faible \\
Qualité des eaux souterraines & De nouvelles sources et de nouveaux modèles de transformation ;
l'architecture est inchangée & moyen \\
Eaux de surface & La chaîne hydrométrique existe déjà ; l'extension porte sur la modélisation &
moyen \\
Autre compartiment (air, sol) & Les sources et les indicateurs diffèrent, le patron tient &
moyen \\
Couplage à un modèle hydrogéologique distribué & Une interface avec un simulateur physique
externe et l'assimilation de ses sorties & élevé \\
\bottomrule
\end{tabularx}
\end{center}

Ce qui est peu coûteux l'est \emph{parce que}
l'infrastructure existe. Le coût élevé de l'année écoulée a été payé une fois ; il ne se
représente pas à chaque extension. C'est la définition même d'un socle.

\section{Ce qui est transmis aux travaux à venir}

Les travaux d'explicabilité se poursuivront au-delà de cette période. Le tableau suivant
confronte ce qu'ils supposent à ce qui leur est effectivement transmis.

\begin{center}
\begin{tabularx}{\textwidth}{@{}p{4.6cm}X@{}}
\toprule
\textbf{Besoin} & \textbf{Disponible} \\
\midrule
Un jeu de données consolidé, joint au forçage climatique & Entrepôt national, partitions depuis
1967 pour les stations et 1950 pour le climat, rattachement spatial matérialisé \\
Des modèles à expliquer & Douze architectures de prévision profonde sous interface unique, avec
traçabilité des expériences \\
Une référence interprétable à laquelle confronter les explications & Modèles à fonction de
transfert calibrés, dont les paramètres ont un sens hydrogéologique \\
Des méthodes d'attribution & Cinq familles complémentaires, confrontables entre elles \\
Des explications actionnables & Contrefactuels sous contrainte physique, mettant en œuvre les
critères formulés par les hydrogéologues \\
Un moyen de valider une explication & Rejeu du scénario dans un modèle métier indépendant \\
Des cas d'étude qualifiés & Détection non supervisée d'influence anthropique, permettant
d'écarter les chroniques impropres \\
\bottomrule
\end{tabularx}
\end{center}

Ce qui n'est pas transmis doit être dit avec la même netteté : aucun résultat de performance
comparée, aucune évaluation quantitative des méthodes contrefactuelles, et aucune campagne
d'expériences consolidée. Le socle rend ces travaux possibles ; il ne les a pas menés.

\section{Ce qui manque pour un véritable environnement d'outils}

Six chantiers
séparent l'état actuel d'un environnement d'outils partagé. Aucun n'est un verrou de recherche ;
tous sont des décisions d'investissement.

\begin{enumerate}
  \item \textbf{L'alignement de l'évapotranspiration sur la chaîne des stations.} C'est le seul
  chantier de la liste qui conditionne la validité de résultats à venir, et non seulement le
  confort de travail. Tant qu'il n'est pas traité, les modèles métier sont calibrés sur un
  forçage deux fois trop élevé, dont le chapitre~\ref{ch:indicateurs} montre qu'il n'est pas
  absorbé par la calibration. Il suppose d'ajouter les températures extrêmes à la chaîne
  station, d'y calculer l'évapotranspiration de référence au pas journalier, de reconstruire
  deux tables partitionnées, puis de recalibrer.

  \item \textbf{L'automatisation des vérifications.} La suite de tests de la plateforme n'est pas
  exécutée automatiquement (chapitre~\ref{ch:apprentissage}). Tant que les vérifications
  dépendent d'une initiative humaine, la confiance dans les évolutions décroît avec le temps.
  C'est le chantier le moins coûteux et le plus rentable.

  \item \textbf{Un accès élargi et gouverné.} L'accès nominatif convient à un outil de projet, non
  à un environnement partagé. Une politique d'accès, précisant qui accède, pour quels usages et avec quelle
  responsabilité sur les données, conditionne tout élargissement.

  \item \textbf{Un catalogue de données.} L'entrepôt est documenté pour qui sait où chercher.
  Un environnement d'outils suppose un catalogue décrivant les jeux disponibles, leur origine,
  leur fraîcheur et leurs conditions d'usage.

  \item \textbf{Le calcul mutualisé.} L'entraînement s'appuie aujourd'hui sur un serveur unique.
  Un usage partagé suppose une file d'attente et un partage de ressources.

  \item \textbf{La campagne d'évaluation comparative.} Le socle sera d'autant plus crédible qu'il
  aura servi à établir un résultat : la comparaison systématique modèles métier / modèles appris,
  et l'évaluation des méthodes contrefactuelles. L'outillage est prêt de part et d'autre.
\end{enumerate}

\section{Articulation avec le volet des jumeaux numériques}

Un jumeau numérique de la ressource en eau suppose trois choses : une représentation de l'état
courant du système, une capacité à le projeter, et une capacité à expliquer la projection. Les
travaux présentés apportent une contribution à chacune.

\begin{center}
\begin{tabularx}{\textwidth}{@{}lX@{}}
\toprule
\textbf{Fonction attendue} & \textbf{Contribution disponible} \\
\midrule
État courant du système & Entrepôt actualisé quotidiennement, indicateurs standardisés validés,
observatoire spatial \\
Projection & Modèles métier calibrés avec scénarisation bornée physiquement, catalogue de
modèles appris \\
Explication de la projection & Attribution multi-méthodes, contrefactuels sous contrainte
physique, détection d'influence anthropique \\
Confiance & Protocole de validation documenté, traçabilité des expériences, limites consignées \\
\bottomrule
\end{tabularx}
\end{center}

La contribution la plus spécifique est la troisième ligne. Un jumeau numérique qui projette sans
expliquer n'est qu'un simulateur ; ce qui en fait un instrument de décision est la capacité à
répondre « pourquoi » et « que se passerait-il si » dans les termes du métier. Les contrefactuels
sous contrainte physique décrits en section~\ref{sec:physcf} sont directement conçus pour cet
usage : ils produisent des scénarios énonçables, du type « un hiver plus sec et plus chaud d'un
degré », plutôt que des perturbations numériques.

Les propriétés qui distinguent une infrastructure d'un prototype sont donc identifiées, et elles
sont effectivement tenues : interface stable, cœur indépendant, indicateurs centralisés,
contrats exécutables, reproductibilité, documentation écrite pour la reprise. Le patron mis en
œuvre se transpose à d'autres compartiments environnementaux.

Six chantiers séparent néanmoins ce socle d'un environnement d'outils réellement partagé, et
aucun d'eux n'est un verrou de recherche. La démonstration la plus convaincante resterait
d'ailleurs qu'il ait servi à établir un résultat scientifique, ce qui n'a pas encore eu lieu.
